\documentclass[11pt]{article}

\usepackage[final]{acl}

\usepackage{times}
\usepackage{latexsym}
\usepackage[T1]{fontenc}
\usepackage[utf8]{inputenc}
\usepackage{microtype}
\usepackage{inconsolata}

\usepackage{amsmath}
\usepackage{booktabs}
\usepackage{graphicx}

\title{Fact-Checking Climate Claims with Sparse Retrieval and Cross-Encoder Reranking}

\author{Group 56}

\begin{document}
\maketitle

\begin{abstract}
We describe a fact-checking system for climate-science claims that retrieves supporting evidence from a 1.2M-passage corpus and classifies each claim as \textsc{Supports}, \textsc{Refutes}, \textsc{Not Enough Info}, or \textsc{Disputed}. The system is built as a modular retrieval-then-classification pipeline: BM25 selects a broad candidate pool, a MiniLM cross-encoder reranks claim-evidence pairs, a relative-delta rule selects a variable number of final evidence passages, and a second MiniLM cross-encoder predicts the claim label from the claim and selected evidence. We tune the pipeline with a sequential greedy ablation rather than a full factorial search, using recall@500 for first-stage retrieval, evidence F-score for selection, and accuracy plus per-class recall for classification. On the development set, the final system obtains evidence F-score 0.2448, label accuracy 0.4805, and harmonic mean 0.3243. The largest gain comes from the retrieval pipeline: final retrieval doubles evidence F-score over raw BM25 top-5 evidence. Classification remains the main weakness, with a strong bias toward \textsc{Supports} and zero recall for \textsc{Refutes}.
\end{abstract}

\section{Introduction}

Automated fact-checking requires both finding relevant evidence and reasoning over it. This coupling is central to FEVER-style verification \citep{thorne2018fever} and is especially important for climate claims, where statements are often paraphrastic, evidence may be multi-sentence, and labels can depend on uncertainty or disagreement rather than simple true/false judgments \citep{diggelmann2020climatefever}. In this assignment, each system receives a claim and must return evidence passage identifiers from a corpus of 1,208,827 passages, together with one of four labels: \textsc{Supports}, \textsc{Refutes}, \textsc{Not Enough Info} (NEI), or \textsc{Disputed}.

The main challenge is that the final score rewards balanced performance. A strong classifier cannot help if the right evidence never enters the candidate pool, while high-recall retrieval is insufficient if the selected passages are noisy or the classifier falls back to majority-class behavior. Our design therefore separates the problem into three questions: can sparse retrieval include the gold evidence in a manageable pool, can a cross-encoder reranker improve the final evidence set, and can a lightweight classifier use that evidence to improve over a majority-label baseline?

We find that the retrieval side is effective but still far from solved. Stemming and BM25 tuning improve recall@500, and MiniLM reranking plus relative-delta evidence selection raises development evidence F-score from 0.1234 for raw BM25 top-5 to 0.2448. The classifier, however, improves accuracy over an all-\textsc{Supports} baseline by only 3.9 percentage points. Error analysis shows that retrieval failures and classification failures are both substantial, and oracle retrieval does not automatically improve label accuracy, suggesting that evidence reasoning remains brittle.

\section{Approach}

\subsection{Task Formulation}

For each claim $c$, the system ranks evidence passages $e \in E$ and returns a selected evidence set $R(c)$ plus a label $\hat{y} \in Y$, where $Y=\{\textsc{Supports},\textsc{Refutes},\textsc{NEI},\textsc{Disputed}\}$. The official score reports evidence retrieval F-score $F$, label accuracy $A$, and their harmonic mean:

\begin{equation}
  H = \frac{2FA}{F + A}.
  \label{eq:harmonic}
\end{equation}

We optimize modules with local development metrics, then evaluate the full pipeline with $F$, $A$, and $H$. This is a sequential greedy ablation: each stage is chosen conditional on earlier stages. This is less exhaustive than a full factorial search, but it keeps training feasible under the assignment's free-Colab constraint.

\subsection{First-Stage Retrieval}

The first stage uses sparse lexical retrieval. We compare TF-IDF and BM25 \citep{robertson2009probabilistic}, with three preprocessing settings: lowercasing only, lowercasing plus stopword removal, and lowercasing plus stopword removal plus Porter stemming \citep{porter1980algorithm}. We choose the first-stage configuration by development recall@500, because evidence that is absent from the top-500 pool cannot be recovered by downstream reranking.

The selected setting is BM25 with lowercasing, stopword removal, Porter stemming, $k_1=1.2$, and $b=0.5$. This retrieves 500 candidates per claim for train, development, and test splits.

\subsection{Cross-Encoder Reranking}

The reranker is cross-encoder/ms-marco-MiniLM-L-12-v2, a MiniLM model \citep{wang2020minilm} trained for MS MARCO passage ranking \citep{bajaj2016msmarco}. Unlike a bi-encoder, the cross-encoder reads the claim and evidence passage jointly and outputs a relevance logit $s(c,e)$.

Training pairs are built from the selected BM25 top-500 pool. Gold claim-evidence pairs are positives; negatives are sampled from BM25 candidates that are not gold evidence. This produces 20,610 training pairs, with 4,122 positives and four hard negatives per positive. We train for up to three epochs and select the checkpoint with the best development evidence F-score under a quick relative-delta check. The first epoch is best; later epochs reduce development F-score, consistent with overfitting to the small pair dataset. Per-epoch numbers are reported in Section~\ref{sec:retrieval-results}.

\subsection{Evidence Selection}

Reranking gives an ordered list and logits, but the submission must return a small final evidence set. We compare fixed-$K$, global-threshold, and relative-delta selection. The final system uses relative-delta selection:

\begin{equation}
  R(c) = \{e_i: s(c,e_i) \geq \max_j s(c,e_j) - \delta\}.
  \label{eq:relative-delta}
\end{equation}

This rule adapts to each claim's score scale instead of assuming that logits are calibrated across claims. The selected threshold is $\delta=1.5$, which yields 4.05 selected evidence passages per claim on average.

\subsection{Claim Classification}

The classifier receives the claim concatenated with up to five selected evidence passages and predicts one of the four labels. During classifier training, the evidence input is the union of gold and retrieved evidence, deduplicated and truncated to five passages; at development and test time, only retrieved evidence is available. We compare distilbert-base-uncased \citep{sanh2019distilbert}, cross-encoder/ms-marco-MiniLM-L-12-v2, and cross-encoder/nli-MiniLM2-L6-H768. Models are trained with learning rate $10^{-5}$ for up to four epochs, batch size 16 for training and 32 for evaluation.

Class imbalance is severe: \textsc{Supports} is the largest class and \textsc{Disputed} the smallest. We therefore compare no class weighting, square-root inverse-frequency weighting, and full inverse-frequency weighting. The final submitted classifier uses MiniLM-MARCO without class weighting, because stronger weighting increases minority-class pressure but sharply damages overall accuracy.

\section{Experiments}

\subsection{Data and Evaluation}

The dataset statistics are shown in Table~\ref{tab:data-summary}. Training labels are imbalanced: 42.26\% \textsc{Supports}, 16.21\% \textsc{Refutes}, 31.43\% \textsc{NEI}, and 10.10\% \textsc{Disputed}. The average claim has 3.357 gold evidence passages, so selecting only one passage is usually insufficient.

\begin{table}[t]
  \centering
  \small
  \begin{tabular}{lr}
    \toprule
    \textbf{Item} & \textbf{Count} \\
    \midrule
    Train claims & 1,228 \\
    Dev claims & 154 \\
    Test claims & 153 \\
    Evidence passages & 1,208,827 \\
    Mean gold evidence / claim & 3.357 \\
    \bottomrule
  \end{tabular}
  \caption{Dataset summary.}
  \label{tab:data-summary}
\end{table}

The official evaluation script computes aggregate evidence F-score by comparing predicted evidence identifiers to gold identifiers, classification accuracy by comparing predicted labels to gold labels, and $H$ from Equation~\ref{eq:harmonic}. We report development results because test labels are unavailable locally; the notebook also generates test-output.json for leaderboard submission.

\subsection{Experimental Setup}

All reported development numbers are produced by v13.ipynb and summarised in v13\_experiment\_summary.md. We use seed 42 for the main run. Reranker and classifier checkpoints are saved under outputs\_notebook\_v13; candidate lists, cross-encoder scores, and final evidence sets are cached for reproducibility. The full search is intentionally sequential: retraining rerankers and classifiers for every preprocessing and BM25 combination would multiply the compute cost.

\subsection{Baselines}

We use two majority-label baselines. Both predict \textsc{Supports} for every claim, matching the largest development label prior. One uses raw BM25 top-5 evidence, and the other uses the final retrieval output. This isolates whether improvements come from evidence retrieval or claim classification.

\section{Results}

\subsection{Retrieval and Selection}
\label{sec:retrieval-results}

Table~\ref{tab:retrieval} summarises the first-stage retrieval comparison. BM25 consistently outperforms TF-IDF, and Porter stemming is the most important preprocessing change. The final first-stage setting reaches 0.6782 recall@500, but its direct top-5 evidence F-score is only 0.1234, motivating reranking and evidence selection.

\begin{table}[t]
  \centering
  \small
  \setlength{\tabcolsep}{3.5pt}
  \begin{tabular}{llrr}
    \toprule
    \textbf{Preproc.} & \textbf{Retriever} & \textbf{R@500} & \textbf{F@5} \\
    \midrule
    stem & BM25 $k_1=1.2,b=0.5$ & 0.6782 & 0.1234 \\
    stem & BM25 $k_1=1.5,b=0.5$ & 0.6758 & 0.1217 \\
    stem & BM25 $k_1=1.2,b=0.75$ & 0.6702 & 0.1266 \\
    stem & TF-IDF & 0.6513 & 0.0919 \\
    stop & BM25 $k_1=1.2,b=0.5$ & 0.6502 & 0.1209 \\
    lower & BM25 $k_1=1.2,b=0.5$ & 0.6194 & 0.1313 \\
    \bottomrule
  \end{tabular}
  \caption{Representative first-stage retrieval results. ``stem'' means lowercasing, stopword removal, and Porter stemming.}
  \label{tab:retrieval}
\end{table}

The reranker peaks after one epoch: training loss decreases from 0.3809 to 0.2632 and 0.2291 over three epochs, but quick development F drops from 0.1961 to 0.1911 and 0.1806. This suggests that the cross-encoder quickly learns useful relevance signals but begins to overfit the 20,610 pair training set.

Selection results in Table~\ref{tab:selection} show that relative-delta selection is better than fixed-$K$ and global thresholds. A global threshold assumes scores are calibrated across claims, but cross-encoder logits vary by claim. The relative rule instead asks how close each evidence passage is to the best passage for the same claim.

\begin{table}[t]
  \centering
  \small
  \begin{tabular}{lrr}
    \toprule
    \textbf{Strategy} & \textbf{Param.} & \textbf{Dev F} \\
    \midrule
    Relative-delta & 1.50 & 0.2448 \\
    Relative-delta & 1.00 & 0.2329 \\
    Relative-delta & 2.00 & 0.2321 \\
    Fixed-$K$ & 3 & 0.2156 \\
    Fixed-$K$ & 4 & 0.2149 \\
    Threshold & 0.90 & 0.2079 \\
    \bottomrule
  \end{tabular}
  \caption{Evidence selection results after reranking.}
  \label{tab:selection}
\end{table}

\subsection{Classification}

Table~\ref{tab:classification} compares classifier backbones with square-root inverse-frequency weighting fixed for this stage. MiniLM-MARCO gives the highest overall accuracy, but the per-class recalls reveal a serious bias: it obtains 0.9706 recall for \textsc{Supports}, 0 for \textsc{Refutes}, and 0 for \textsc{Disputed}. The NLI-oriented MiniLM is less accurate overall but recovers some \textsc{Refutes} and \textsc{Disputed} examples, suggesting that the final classifier's accuracy is partly driven by label-prior behavior rather than robust evidence reasoning.

\begin{table}[t]
  \centering
  \small
  \setlength{\tabcolsep}{3pt}
  \begin{tabular}{lrrrr}
    \toprule
    \textbf{Backbone} & \textbf{A} & \textbf{Supp.} & \textbf{Ref.} & \textbf{Disp.} \\
    \midrule
    MiniLM-MARCO & 0.4935 & 0.9706 & 0.0000 & 0.0000 \\
    MiniLM-NLI & 0.4805 & 0.8088 & 0.3333 & 0.1111 \\
    DistilBERT & 0.4481 & 0.8088 & 0.0370 & 0.0000 \\
    \bottomrule
  \end{tabular}
  \caption{Classifier backbone comparison. Columns after accuracy show per-class recall for selected classes.}
  \label{tab:classification}
\end{table}

After selecting MiniLM-MARCO, we compare class weighting in Table~\ref{tab:weighting}. Full inverse-frequency weighting raises \textsc{Disputed} recall to 0.8889 but collapses accuracy to 0.2468. Square-root weighting gives the best $H$ by a narrow margin, but predicts no \textsc{Disputed} claims correctly. We choose no weighting for the final system because it preserves some \textsc{Disputed} recall while avoiding the large accuracy collapse.

\begin{table}[t]
  \centering
  \small
  \setlength{\tabcolsep}{4pt}
  \begin{tabular}{lrrrr}
    \toprule
    \textbf{Weighting} & \textbf{A} & \textbf{H} & \textbf{Ref.} & \textbf{Disp.} \\
    \midrule
    None & 0.4805 & 0.3243 & 0.0000 & 0.2778 \\
    Sqrt inv. freq. & 0.4870 & 0.3258 & 0.0000 & 0.0000 \\
    Inv. freq. & 0.2468 & 0.2458 & 0.0741 & 0.8889 \\
    \bottomrule
  \end{tabular}
  \caption{Class weighting ablation for the MiniLM-MARCO classifier.}
  \label{tab:weighting}
\end{table}

\subsection{Final Development Results}

Table~\ref{tab:main-results} compares the final system to majority-label baselines. The final retrieval output doubles evidence F-score over BM25 top-5 evidence, improving $H$ from 0.1929 to 0.3150 even when the label is always \textsc{Supports}. The classifier then raises accuracy from 0.4416 to 0.4805, giving final $H=0.3243$. Thus most of the measured improvement comes from retrieval, while classification contributes a smaller but real gain.

\begin{table}[t]
  \centering
  \small
  \setlength{\tabcolsep}{4pt}
  \begin{tabular}{lrrr}
    \toprule
    \textbf{System} & \textbf{F} & \textbf{A} & \textbf{H} \\
    \midrule
    All \textsc{Supports} + BM25 top-5 & 0.1234 & 0.4416 & 0.1929 \\
    All \textsc{Supports} + final retrieval & 0.2448 & 0.4416 & 0.3150 \\
    Final system & 0.2448 & 0.4805 & 0.3243 \\
    \bottomrule
  \end{tabular}
  \caption{Main development-set results.}
  \label{tab:main-results}
\end{table}

The confusion matrix in Table~\ref{tab:confusion} makes the classifier limitation clear. The system predicts \textsc{Supports} for most claims and never predicts \textsc{Refutes} correctly. This is the main reason accuracy remains close to the majority baseline despite better evidence retrieval.

\begin{table}[t]
  \centering
  \small
  \setlength{\tabcolsep}{3pt}
  \begin{tabular}{lrrrr}
    \toprule
    \textbf{Gold} & \textbf{Sup.} & \textbf{Ref.} & \textbf{NEI} & \textbf{Disp.} \\
    \midrule
    \textsc{Supports} & 65 & 0 & 2 & 1 \\
    \textsc{Refutes} & 24 & 0 & 3 & 0 \\
    \textsc{NEI} & 37 & 0 & 4 & 0 \\
    \textsc{Disputed} & 13 & 0 & 0 & 5 \\
    \bottomrule
  \end{tabular}
  \caption{Final classifier confusion matrix on the development set. Rows are gold labels and columns are predictions.}
  \label{tab:confusion}
\end{table}

\subsection{Error Analysis}

Seed variance confirms that the classifier result should not be overinterpreted. With the same retrieval fixed and seeds 42, 123, and 2024, mean accuracy is $0.4892 \pm 0.0201$ and mean $H$ is $0.3262 \pm 0.0046$. Retrieval F-score is unchanged because evidence selection is fixed.

Because \textsc{Disputed} claims often require detecting conflict between evidence passages, we also tried a lightweight Qwen2.5-3B-Instruct cascade. For claims currently predicted as \textsc{Supports} or \textsc{Refutes} with at least two retrieved passages, the model was prompted to answer whether the passages contradicted each other; a positive answer would change the label to \textsc{Disputed}. As Table~\ref{tab:llm-cascade} shows, the cascade made no label changes on the 130 eligible development claims, so it did not improve the final metrics. This negative result suggests that a zero-shot contradiction prompt is not enough; a useful conflict module would likely need stronger evidence-level NLI supervision or calibration.

\begin{table}[t]
  \centering
  \small
  \begin{tabular}{lrrr}
    \toprule
    \textbf{Setting} & \textbf{F} & \textbf{A} & \textbf{H} \\
    \midrule
    Final system & 0.2448 & 0.4805 & 0.3243 \\
    LLM cascade & 0.2448 & 0.4805 & 0.3243 \\
    \bottomrule
  \end{tabular}
  \caption{Qwen2.5-3B-Instruct \textsc{Disputed} cascade on 130 eligible development claims. The cascade flipped zero labels.}
  \label{tab:llm-cascade}
\end{table}

We also run an oracle diagnostic by replacing retrieved development evidence with gold evidence. This raises evidence F-score to 1.0 and therefore $H$ to 0.5806, but label accuracy falls from 0.4805 to 0.4091. The drop implies that the classifier is not simply waiting for cleaner evidence; it is sensitive to the evidence distribution seen at training and may rely on label priors or noisy retrieval patterns.

Finally, we bucket errors by whether retrieval hits at least one gold evidence passage and whether the label is correct. Only 27.9\% of claims are both retrieval-hit and label-correct. Another 24.7\% have a retrieval hit but a wrong label, 27.3\% fail both retrieval and classification, and the remaining 20.1\% are ``lucky'' cases where the label is correct without any gold evidence in the retrieved set, most likely from label-prior or claim-text shortcuts. This split shows that future improvements must target both candidate recall and evidence-conditioned reasoning.

\section{Conclusion}

Our final system is a compact, reproducible retrieval-then-classification pipeline for climate claim verification. The strongest component is retrieval: BM25 tuning, cross-encoder reranking, and relative-delta evidence selection substantially improve evidence F-score over raw BM25 top-5. The main limitation is classification. The MiniLM-MARCO classifier slightly improves over the majority-label baseline, but its strong \textsc{Supports} bias and zero \textsc{Refutes} recall show that it has not learned robust stance reasoning.

The most promising future directions are therefore targeted rather than architectural: iterative hard-negative mining for the reranker, per-evidence NLI scoring followed by learned aggregation, and training objectives that improve minority-class recall without collapsing overall accuracy. A fuller hyperparameter search could also expose interactions missed by the sequential greedy ablation, but the current design keeps the submitted notebook feasible under the compute constraints.

\section*{Team Contributions}

\begin{itemize}
  \item Mengfan Gao (1660072): contributed to classifier experiments, ablation analysis, error analysis, and report writing.
  \item Jingwang Liu (1442608): contributed to data processing, baseline comparisons, reproducibility checks, and report editing.
  \item Kexing Ma (1697372): contributed to retrieval experiments, reranker development, evaluation, and report writing.
\end{itemize}

\bibliography{custom}

\end{document}
